\documentclass[10pt,a4paper]{article}
\usepackage{fontspec,amsmath,amssymb,graphicx}
\usepackage[margin=18mm]{geometry}
\setmainfont{Times New Roman}
\setlength{\parindent}{0pt}
\setlength{\parskip}{5pt}
\allowdisplaybreaks
\sloppy
\emergencystretch=2em
\begin{document}
\section*{VD\_4\_1\_2\_03 — Lời giải bổ sung}
\subsection*{Kiến thức và trình tự giải}
\textbf{Bước 1 — Chọn đúng đại lượng.} Năng lượng bình phương là $I(t)=\frac12\int_0^Lu^2dx$. Để chứng minh duy nhất, áp dụng cho hiệu $w$ của hai nghiệm cùng dữ kiện; $w$ có dữ liệu đầu bằng không và biên thuần nhất. Tổng nhiệt $E=\int_0^Lu\,dx$ là đại lượng khác, dùng khi đề hỏi bảo toàn hoặc suy giảm tổng nhiệt.

\textbf{Bước 2 — Đạo hàm và tích phân từng phần.} Với $u_t=Ku_{xx}$,


\[
I'=\int_0^Luu_t\,dx
=K[uu_x]_0^L-K\int_0^Lu_x^2dx.
\]


Thay biên tại $L$ rồi trừ giá trị tại $0$. Dirichlet/Neumann thuần nhất làm hạng biên bằng không. Robin $u_x(0)=u(0)$ đóng góp $-u(0)^2$; Robin $u_x(L)=-u(L)$ đóng góp $-u(L)^2$. Biên tuần hoàn triệt tiêu hai đầu. Một điều kiện ghép khác phải được kiểm tra trực tiếp, vì hạng biên có thể dương.

\textbf{Bước 3 — Suy duy nhất và tính dương của trị riêng.} Nếu $I'\le0$ và $I(0)=0$ thì $0\le I(t)\le0$, nên $w=0$. Với $X''+\lambda X=0$, nhân với $X$ rồi tích phân:


\[
\lambda\int_0^LX^2dx=\int_0^L(X')^2dx-[XX']_0^L.
\]


Nếu vế phải không âm thì $\lambda\ge0$. Khi bằng không, buộc từng số hạng không âm bằng không và kiểm tra biên để quyết định có mode hằng hay chỉ nghiệm tầm thường.

\textbf{Bước 4 — Chứng minh trực giao.} Nhân hai phương trình riêng lần lượt với $X_m,X_n$, trừ rồi tích phân:


\[
(\lambda_n-\lambda_m)\int_0^LX_nX_mdx
=[X_nX_m'-X_n'X_m]_0^L.
\]


Biên tự liên hợp làm vế phải bằng không; $\lambda_n\ne\lambda_m$ suy trực giao. Với biên không tự liên hợp hoặc đạo hàm bậc ba, không tự động dùng công thức chiếu trực giao.

\textbf{Bước 5 — Khi có hạng phản ứng hoặc bậc đạo hàm khác.} Với Dirichlet và $u_t=Ku_{xx}+Au_x+Bu$, hạng $A\int uu_x=A[u^2]_0^L/2=0$, nên $I'=-K\int u_x^2+2BI\le2BI$; $e^{-2Bt}I$ không tăng, vẫn suy duy nhất. Với $u_t=Ku_{xxx}$, tích phân hai lần cho $I'=K[uu_{xx}-u_x^2/2]_0^L$. Với tổng nhiệt trong bài Neumann có $-ru$, $E'=K[u_x]_0^L-rE$; giải ODE này trước khi kết luận.
\subsection*{Lời giải áp dụng cho bài}
\textbf{Bước 1. Năng lượng và duy nhất.} Biên là $u_x(0)=u(0)$, $u_x(L)=0$. Tích phân từng phần:
\[I'=\int uu_t=K[uu_x]_0^L-K\int u_x^2
=-K u(0,t)^2-K\int_0^Lu_x^2dx\le0.\]
Áp dụng cho hiệu hai nghiệm có dữ liệu đầu $0$ cho $I(0)=0$, nên hiệu bằng không.

\textbf{Bước 2. Trị riêng và trực giao.} Thế $u=XT$ cho $X''+\lambda X=0$, $X'(0)=X(0)$, $X'(L)=0$, $T'=-K\lambda T$. Nhân ODE với $X$:
\[\lambda\int X^2=\int(X')^2+X(0)^2.\]
Nếu $\lambda=0$, $X'=0$ và $X(0)=0$, suy $X=0$; mọi trị riêng không tầm thường có $\lambda>0$. Trừ hai phương trình như đẳng thức Green cho $(\lambda_n-\lambda_m)\int X_nX_m=[X_nX_m'-X_n'X_m]_0^L=0$, nên trực giao.

\textbf{Bước 3. Tìm $k$ và chuẩn hóa.} Với $\lambda=k^2$, $X=A\cos kx+B\sin kx$. Biên trái cho $A=kB$; chọn $A=1$, được $X=\cos kx+\sin kx/k$. Biên phải cho $-k\sin kL+\cos kL=0$, tức $k\tan kL=1$. Trên mỗi khoảng $((n-1)\pi/L,(n-1/2)\pi/L)$, $k\tan kL$ tăng từ $0$ đến $\infty$, nên có đúng một nghiệm $k_n$, $n\ge1$. Chuẩn là
\[N(k)=\frac L2(1+k^{-2})+\frac{\sin2kL}{4k}(1-k^{-2})+\frac{\sin^2kL}{k^2}.\]
Do đó $u=\sum c_ne^{-Kk_n^2t}X_n$, $c_n=\int fX_n/N(k_n)$.

\textbf{Bước 4. Tính dữ liệu cụ thể $K=2,L=\pi,f=1+2\sin(x/2)$.}
Đặt $q=1/2$, và với $r\ne0$ đặt $D(r)=(1-\cos\pi r)/r$, $E(r)=\sin\pi r/r$, mở rộng $D(0)=0,E(0)=\pi$ theo giới hạn. Tích–tổng và tích phân từng hạng cho
\begin{align*}
J(k)=\int_0^\pi fX\,dx
={}&\frac{\sin\pi k}{k}+\frac{1-\cos\pi k}{k^2}\\
&+D(q+k)+D(q-k)+\frac{E(q-k)-E(q+k)}k.
\end{align*}
Vậy $c_n=J(k_n)/N(k_n)$ và
\[u(x,t)=\sum_{n\ge1}\frac{J(k_n)}{N(k_n)}e^{-2k_n^2t}
\left(\cos k_nx+\frac{\sin k_nx}{k_n}\right).\]
Không còn tích phân hệ số chưa tính. Với $L=\pi$, các nghiệm đầu là $0.3834486028,1.2187220137,2.1391968186$. Mỗi hạng thỏa hai biên và PDE. $f'(0)=1=f(0)$, $f'(\pi)=0$, nên dữ liệu khớp biên; chuỗi trả lại $f$ theo khai triển Sturm–Liouville và trơn khi $t>0$. Tất cả $k_n>0$, suy giới hạn đều bằng $0$. Không mặc định các đạo hàm cấp cao liên tục ở góc, vì điều đó cần thêm tương thích.
\begin{center}\includegraphics[width=.97\linewidth,height=11cm,keepaspectratio]{Do_thi/VD_4_1_2_03.png}\end{center}
\textit{Đồ thị minh họa; với nghiệm số, xem phương pháp và tham số trong mã Python kèm theo.}\end{document}